\documentclass[10.5pt,a4paper]{article}
\usepackage[T1]{fontenc}
\usepackage[utf8]{inputenc}
\usepackage{lmodern}
\usepackage{textcomp}
\usepackage{amsmath}
\usepackage{placeins}
\usepackage[margin=2.3cm]{geometry}
\usepackage{microtype}
\usepackage{graphicx}
\usepackage{booktabs}
\usepackage{xcolor}
\usepackage{caption}
\usepackage{enumitem}
\usepackage[hidelinks]{hyperref}
\captionsetup{font=small,labelfont=bf}
\setlist{itemsep=1pt,topsep=3pt}
\setlength{\parskip}{4pt}
\setlength{\parindent}{0pt}
\setlength{\emergencystretch}{3em}
\renewcommand{\topfraction}{0.9}
\renewcommand{\textfraction}{0.1}
\newcommand{\code}[1]{\texttt{#1}}
\definecolor{good}{HTML}{1c5cab}
\newcommand{\best}[1]{\textbf{\textcolor{good}{#1}}}

\title{\textbf{Learned Page Replacement in xv6}\\[2pt]
\large Week 15 update: new workloads, a wider model search, and the kernel}
\author{}
\date{}

\begin{document}
\maketitle
\vspace{-2em}

\section*{Summary}

Week 14 showed that a small linear model over signals the kernel can see (accessed and dirty bits and cheap
bookkeeping) beats Clock by 10--15\% in faults on our benchmarks. This week tested how far that result
reaches. All numbers are faults relative to Clock (below 1 is better), averaged geometrically over memory
sizes of 5, 10 and 20\% of each program's pages.

\begin{itemize}
\item \textbf{More, harder workloads.} Added the fixed \code{lzw} benchmark to the full pipeline, a synthetic
  pattern suite with known answers (\code{patbench}), replays of real Twitter cache traces, real Redis and
  SQLite page traces, and four application-shaped programs (pointer chase, sparse matrix--vector product, hash
  join, Bloom filter). Every program's access string is checked reference by reference against an independent
  model.
\item \textbf{Literature baselines.} ARC, S3-FIFO, an MGLRU analogue, LHD and exact MRU are implemented in the
  simulator and checked against independent reference implementations. LHD is the strongest untrained policy
  (0.90--0.93 on our data); ARC and S3-FIFO are mixed, and S3-FIFO is worse than Clock on graph, matmul and sort.
\item \textbf{Model zoo.} Trees (CART, random forest, boosted trees, ranking boosted trees, logistic regression)
  and neural nets (deeper MLP, set-attention, CNN and LSTM over page history, a CNN over the kernel's own
  scan-bit history, mixed models) were fitted on every scope and selected on validation only. With signals the
  kernel can actually observe, \best{no model class beats the linear scorer by the 2\% and at-most-2$\times$-cost
  bar overall}. Gains of 5--12\% appear on single workloads (kv, the Twitter replays);
  models given oracle information gain up to 6\% in distribution and fail badly on the real Redis and SQLite traces.
\item \textbf{In the xv6 kernel} (157 runs on the new workloads, plus 61 with sampled scoring): the learned
  policy cuts faults by 16\% on lzw (0.84), 15\% on the Twitter replays (0.85) and 15\% on patbench (0.85)
  against Clock, and does nothing on the application programs (1.00). The simulator predicts kernel faults to
  within about 4\%, with one unexplained exception (Section~\ref{sec:kernel}). Sampled scoring, now in the kernel, roughly halves the selection cost for 1--2.5\% more faults.
\item \textbf{A reproducibility bug was found and fixed.} Training samples depended on Python's per-process
  string hashing, so datasets could not be rebuilt identically. All models were refit; Clock, LRU and
  the other classical results are unchanged.
\end{itemize}

\section{What changed in the data and the pipeline}

\textbf{Reproducibility fix.} The sampling seed of each training stream used Python's \code{hash()}, which is
randomised per process, so two builds of the same dataset differed. The seed is now a CRC32 of the stream name, a
capacity and an optional salt. Rebuilding twice is bit-identical (checked over 11 dataset files); all 4\,272
classical result rows of the original streams are unchanged. Because the old caches could not be recovered,
every model was refit, so numbers move slightly from week 14 (for example the global default kernel model now
has five features instead of four). The noise floor from sampling alone is measured in Section~\ref{sec:noise}.

\textbf{lzw.} The corrected lzw benchmark from the team (\code{lzw-r20} for training, validation and test;
\code{lzw-r30} held out) now runs through the whole pipeline: simulator sweep, neural models, selection, test,
kernel export.

\section{New workloads}

\begin{table}[h]
\centering\small
\begin{tabular}{p{2.4cm}p{6.8cm}p{4.6cm}}
\toprule
workload & what it does & why it is in the study \\
\midrule
\code{patbench} & 11 synthetic modes over 1\,024 pages: cyclic loop, scan with a hot set (two sizes),
Zipf with skew 0.6--1.2, uniform, hot set that moves (two speeds), loop/Zipf switching & Known answers: a loop defeats LRU, FIFO and Clock and
is solved by MRU; \code{uniform} has nothing to learn, so any gain there is a leak \\
\code{kv-T} (Twitter) & kvbench driven by real key streams from six production cache clusters, keys hash-sampled
to kvbench's 1\,024-key space (4\,700--40\,000 operations) & Real popularity and write mix; whole
clusters are held out \\
real traces & Redis and SQLite page-touch strings collected on a host (12.3M and 8.6M references, no read/write
marks) & Programs we did not write; tests transfer to a different kind of recording \\
\code{appbench} & pointer chase (mcf-like), CSR sparse matrix--vector, hash join, Bloom filter, each with an
explicit multi-array layout & Irregular, scan-plus-random structure the other programs lack \\
\bottomrule
\end{tabular}
\caption{Workloads added this week. All are held out of training: they only test.}
\end{table}

\textbf{How strings are validated.} Each program's reference string has an exact host-side model. The kernel's
trace equals the model line for line for all 11 patbench modes (24 short runs and 9 full runs), and for all
four appbench modes (30\,000 references each, with matching read, write and checksum counters). The remaining
production traces were then generated from the model rather than by booting xv6 for 17--45 minutes each; their logs carry
a marker saying so and no kernel fault counters. The Twitter streams and the real traces are not
generated: kvbench ran them in the kernel, and the Redis/SQLite files are the host recordings.

\begin{table}[h]
\centering\small
\begin{tabular}{lrrrr}
\toprule
workload & variants & streams & references & held-out streams \\
\midrule
app & 4 & 12 & 12.0M & 12 \\
btree & 5 & 25 & 7.1M & 5 \\
graph & 4 & 14 & 77.9M & 3 \\
kv & 13 & 41 & 8.5M & 11 \\
lzw & 2 & 12 & 61.4M & 6 \\
matmul & 6 & 6 & 13.8M & 2 \\
pat & 11 & 55 & 55.0M & 10 \\
redis & 1 & 1 & 12.3M & 1 \\
sort & 3 & 9 & 16.6M & 3 \\
sqlite & 1 & 1 & 8.6M & 1 \\
\bottomrule
\end{tabular}

\caption{The study dataset after this week (176 streams).}
\end{table}

\section{Literature baselines}

Implemented in the simulator and cross-checked against independent reference implementations on 475 cases
(identical counts): ARC (integer adaptation), S3-FIFO (as published, and a form that reads accessed bits only
at queue heads), an MGLRU analogue for anonymous memory (generations, aging walk, rescue at eviction), LHD in
a scan form, and exact MRU. MRU equals Belady on cyclic loops (9 of 9 checks), and LRU, FIFO and Clock miss every
access of a loop larger than memory, as theory says. Table~\ref{tab:zoo} shows them next to the learned models.

\section{Transfer to unseen workloads}

\begin{figure}[h]
\centering\includegraphics[width=\linewidth]{figs/transfer.pdf}
\caption{Faults relative to Clock on workloads no model was trained on. Linear (K+) is the global linear model
over all kernel signals.}
\end{figure}

\begin{itemize}
\item \best{patbench known answers hold.} On \code{loop} FIFO, LRU, Clock, ARC and LHD miss everything and MRU
  matches Belady; on \code{uniform} every online policy is within 0.0003 of $1-C/P$, so \code{uniform} is a clean
  leak control: no model gains there.
\item The linear kernel-signal model beats ARC, S3-FIFO and LHD on the skewed and scan-plus-hot patterns
  (for example 0.76 against 0.85, 0.81 and 0.90 on Zipf 0.99) and on the Twitter replays (in the kernel, Section~\ref{sec:kernel}).
\item \textbf{Where it loses.} Moving-hot-set patterns, SpMV and hash join. The oracle-feature variant (F+K)
  raises faults 2.6-fold on the moving hot set and reaches the 3$\times$-Clock cap on Redis. The global kernel-signal linear model is 24\% and 14\%
  worse than Clock on SpMV and hash join. Real Redis/SQLite traces carry no write information, so dirty-bit
  signals are absent there, and LRU (0.93, 0.91) is the best simple policy.
\end{itemize}

\begin{table}[h]
\centering\scriptsize
\begin{tabular}{lrrrrrrrr}
\toprule
stream & LRU & ARC & S3-FIFO & LHD & Belady & linear K+ & linear F+K & linear kv K+ \\
\midrule
app-bloom & 1.000 & 0.989 & 0.962 & 0.998 & 0.635 & 0.930 & 0.926 & -- \\
app-chase & 0.991 & 1.002 & 1.021 & 0.984 & 0.683 & 1.021 & 1.020 & -- \\
app-hashjoin & 0.998 & 1.000 & 1.012 & 0.994 & 0.544 & 1.138 & 1.017 & -- \\
app-spmv & 0.992 & 1.000 & 1.026 & 1.046 & 0.727 & 1.239 & 0.995 & -- \\
kv-T & 0.983 & 0.983 & 1.028 & 0.916 & 0.611 & -- & -- & 0.980 \\
pat-loop & 1.000 & 1.000 & 1.000 & 1.000 & 0.882 & 0.894 & 0.887 & -- \\
pat-phase & 0.968 & 0.887 & 0.899 & 0.876 & 0.552 & 0.899 & 2.644 & -- \\
pat-phaseshort & 0.967 & 0.888 & 0.909 & 0.879 & 0.552 & 0.918 & 2.686 & -- \\
pat-scanhot & 0.992 & 0.786 & 0.795 & 0.866 & 0.700 & 0.748 & 0.744 & -- \\
pat-scanhotlo & 1.001 & 0.870 & 0.874 & 0.982 & 0.803 & 0.829 & 0.903 & -- \\
pat-switch & 1.005 & 0.882 & 0.877 & 1.007 & 0.704 & 0.817 & 0.805 & -- \\
pat-uniform & 1.000 & 1.000 & 1.000 & 1.000 & 0.658 & 1.000 & 1.000 & -- \\
pat-zipf060 & 0.998 & 0.915 & 0.881 & 0.975 & 0.612 & 0.845 & 0.824 & -- \\
pat-zipf080 & 0.994 & 0.874 & 0.837 & 0.932 & 0.587 & 0.796 & 0.773 & -- \\
pat-zipf099 & 0.990 & 0.846 & 0.806 & 0.897 & 0.565 & 0.764 & 0.737 & -- \\
pat-zipf120 & 0.984 & 0.825 & 0.783 & 0.863 & 0.544 & 0.737 & 0.710 & -- \\
redis-real & 0.929 & 0.967 & 1.104 & 0.973 & 0.667 & 1.106 & 3.000 & -- \\
sqlite-real & 0.908 & 0.947 & 1.082 & 0.927 & 0.567 & 0.934 & 1.318 & -- \\
\bottomrule
\end{tabular}

\caption{Faults / Clock per stream group (geometric mean over streams and the three memory sizes). Simulator results; the three last columns are linear models (global kernel signals, global oracle+kernel signals, kv-scope kernel signals).}
\end{table}

\section{Model zoo}

\textbf{Protocol.} The grid was written before any result: kernel-observable inputs (K) or all inputs (FK),
with or without candidate-set rank features; CART (depth 4, 8), random forest, boosted trees (30 and 100 trees,
8 and 16 leaves), ranking boosted trees, logistic regression; eight neural architectures. Candidate
sets come from the same recorded eviction decisions as before; the target is the log of the next-use distance.
Probation is fixed at 2 scans. Selection is on validation streams at 10\% memory only; the chosen models are
then replayed on test, held-out and transfer streams at 5, 10 and 20\%. Intervals are 95\% stream-level
bootstraps. The neural models run as numpy callbacks in the simulator (checked against PyTorch to $10^{-6}$)
and are not run on streams of more than 400\,000 references for cost; the C tree scorer is checked against
scikit-learn and LightGBM to $3\times10^{-6}$. A follow-up added four- to sixteen-tree boosted ensembles to
meet the cost bar.

\begin{figure}[h]
\centering\includegraphics[width=\linewidth]{figs/zoo.pdf}
\caption{Faults relative to the best linear scorer (below 1: the model is better), with 95\% intervals.
Models with oracle inputs (blue-grey, orange-red, ``oracle'') cannot run in a kernel.}
\label{fig:zoo}
\end{figure}

\begin{table}[h]
\centering\scriptsize
\begin{tabular}{llrrrrrrrrrrr}
\toprule
scope & split & LRU & ARC & S3F & LHD & linear K & tree K & NN K & NN K$^{\dagger}$ & tree FK & NN FK & Belady \\
\midrule
btree & test & 0.994 & 0.950 & 0.862 & 0.948 & 0.835 & 0.834 & 0.822 & 0.817 & 0.829 & 0.811 & 0.550 \\
btree & held-out & 0.997 & 0.975 & 0.908 & 0.974 & 0.875 & 0.867 & -- & 0.851 & 0.853 & 0.847 & 0.615 \\
graph & test & 0.990 & 0.991 & 1.161 & 0.865 & 0.847 & 0.847 & -- & 0.844 & 0.806 & 0.844 & 0.456 \\
graph & held-out & 0.992 & 0.993 & 1.256 & 0.816 & 0.798 & 0.765 & -- & 0.760 & 0.742 & 0.760 & 0.482 \\
kv & test & 0.984 & 0.986 & 1.047 & 0.906 & 0.897 & 0.902 & 0.850 & 0.849 & 0.875 & 0.843 & 0.559 \\
kv & held-out & 0.982 & 0.982 & 1.013 & 0.901 & 0.902 & 0.903 & 0.845 & 0.858 & 0.900 & 0.854 & 0.577 \\
lzw & test & 0.997 & 0.968 & 0.889 & 0.892 & 0.821 & 0.933 & -- & 0.809 & 0.933 & 0.803 & 0.466 \\
lzw & held-out & 0.981 & 0.957 & 0.866 & 0.887 & 0.813 & 1.068 & -- & 0.841 & 1.068 & 0.812 & 0.409 \\
matmul & held-out & 0.879 & 0.875 & 1.508 & 1.036 & 0.671 & 0.803 & 0.868 & 0.868 & 0.899 & 0.868 & 0.617 \\
sort & test & 1.012 & 1.013 & 1.321 & 1.008 & 0.983 & 0.879 & -- & 0.891 & 0.854 & 0.901 & 0.830 \\
sort & held-out & 0.999 & 0.999 & 1.084 & 1.019 & 0.933 & 0.928 & -- & 0.961 & 0.938 & 0.924 & 0.896 \\
global & test & 0.988 & 0.970 & 0.960 & 0.925 & 0.870 & 0.890 & 0.889 & 0.850 & 0.842 & 0.841 & 0.555 \\
global & held-out & 0.982 & 0.982 & 1.013 & 0.901 & 0.891 & 0.898 & 0.930 & 0.910 & 0.885 & 0.886 & 0.577 \\
\bottomrule
\end{tabular}

\caption{Faults / Clock on test and held-out streams. ``linear K'' is the week-14 model class (kernel signals);
``tree K'' and ``NN K'' use kernel signals only; ``NN K$^{\dagger}$'' also reads the oracle reuse history; ``FK''
adds oracle inputs. A dash: not run (cost rule) or not selected.}
\label{tab:zoo}
\end{table}

\textbf{What the zoo says.}
\begin{itemize}
\item \best{Richer models help only where they are given information a kernel does not have.} With oracle inputs
  the best tree and neural models are 2--5\% better than the best linear model in distribution; with kernel
  inputs only, trees tie the linear model (btree 0.834 against 0.835, graph held-out 0.765 against 0.798,
  global 0.890 against 0.870) and small boosted ensembles (4--16 trees) are \emph{worse}: lzw 1.19 and
  1.57 against 0.82 and 0.81 for the linear model.
\item The deployable neural nets (deep MLP, set-attention, aging-bit CNN) gain 5--6\% on kv and 2\% on btree
  and lose 2--4\% globally. Their kv gain is the one result that passes the 2\% bar on a workload scope, at roughly
  100$\times$ the arithmetic of a linear scorer.
\item \textbf{Robustness is the real difference.} On the Twitter and application streams the kernel-input trees
  beat the linear model on the Twitter replays by 12\% (0.86 against 0.98) and on the application streams by 3.5\%
  (1.04 against 1.08), because they do not extrapolate the way a linear score does; but on the real Redis and SQLite traces every complex model is 22--41\% worse than the linear one.
  None of the extra capacity is free of this trade.
\item matmul has no validation streams (selection used its training streams), and every richer model is
  20--30\% worse than the linear one on its held-out streams: more capacity overfits when there is little data.
\end{itemize}

\section{In the xv6 kernel}
\label{sec:kernel}

The refit linear models were re-exported to integer weights and the learned policy ran in the real kernel on
the new workloads (patbench and appbench for 50\,000 accesses; lzw and the Twitter replays in full) at 10\%
memory, all policies, same frame counts as the simulator.

\begin{figure}[h]
\centering\includegraphics[width=\linewidth]{figs/kernel.pdf}
\caption{Kernel faults relative to Clock; the number on each bar is the selection cost in timer ticks
per eviction (Clock: 24). ``@k'' scans and scores only k random candidates per eviction.}
\end{figure}

\begin{table}[h]
\centering\small
\begin{tabular}{llrrr}
\toprule
workload & policy & runs & faults / Clock & select ticks / eviction \\
\midrule
lzw & fifo & 7 & 1.132 & 8 \\
lzw & clock & 7 & 1.000 & 25 \\
lzw & aging & 7 & 1.050 & 90 \\
lzw & lfu & 7 & 1.060 & 90 \\
lzw & ml global & 7 & 0.941 & 202 \\
lzw & ml lzw & 7 & 0.837 & 174 \\
lzw & ml lzw@64 & 7 & 0.837 & 173 \\
lzw & ml lzw@32 & 7 & 0.842 & 136 \\
lzw & ml lzw@16 & 7 & 0.858 & 93 \\
kv-T & fifo & 6 & 1.315 & 8 \\
kv-T & clock & 6 & 1.000 & 24 \\
kv-T & aging & 6 & 1.123 & 104 \\
kv-T & lfu & 6 & 1.180 & 102 \\
kv-T & ml global & 6 & 0.846 & 223 \\
kv-T & ml global@32 & 6 & 0.849 & 157 \\
kv-T & ml kv & 6 & 0.876 & 263 \\
kv-T & ml kv@32 & 6 & 0.875 & 153 \\
kv-T & ml kv@16 & 6 & 0.884 & 104 \\
pat & fifo & 7 & 1.195 & 10 \\
pat & clock & 7 & 1.000 & 22 \\
pat & aging & 7 & 1.143 & 158 \\
pat & lfu & 7 & 1.148 & 159 \\
pat & ml global & 7 & 0.853 & 428 \\
pat & ml global@32 & 7 & 0.861 & 182 \\
pat & ml global@16 & 7 & 0.873 & 117 \\
pat & ml global-k & 7 & 0.950 & 370 \\
app & fifo & 4 & 1.221 & 10 \\
app & clock & 4 & 1.000 & 24 \\
app & aging & 4 & 0.983 & 133 \\
app & lfu & 4 & 1.092 & 134 \\
app & ml global & 4 & 1.004 & 308 \\
app & ml global@32 & 4 & 0.996 & 148 \\
app & ml global@16 & 4 & 0.992 & 103 \\
app & ml global-k & 4 & 0.997 & 273 \\
\bottomrule
\end{tabular}

\caption{In-kernel runs (faults = zero-fill plus swap faults; geometric means over streams).}
\end{table}

\begin{itemize}
\item \best{The simulator predicts the kernel.} Kernel over simulator faults, geometric mean per policy:
  0.98 on lzw, 0.96--0.98 on patbench and appbench (prefix replay), and 0.98--1.01 for the Twitter replays
  with the global and kernel-only models. The differences come from the kernel's own pages (code, stack) and
  capacity cliffs.
\item \textbf{One exception, not explained.} For the kv-scope model on the Twitter replays the kernel is about 10\% better
  than the simulator (0.90; kernel 0.88 of Clock against 0.98 in the simulator); the same model agrees to 0.3\%
  on the other kv streams in week 14. Clock, the global model and the kernel-only kv model agree on the
  replays. A candidate cause is that the replay program reads its key file through a small buffer outside the arena,
  which the simulator's trace does not contain; it has not been tested. The Twitter-replay kernel numbers for the kv model
  should be read with that in mind.
\item The learned policy saves 16\% of faults on lzw (per-workload model; 6\% for the global model), 15\% on the
  Twitter replays and 15\% on patbench, and does not help on the application programs.
\item \textbf{Cost, and sampled scoring.} The full scorer costs 170--430 ticks per eviction against 24 for
  Clock (7--17$\times$). Scanning and scoring a random sample of 32 candidates takes 40--80\% of that (136
  ticks on lzw, 157 on the Twitter replays, 182 on patbench) for 0--1\% more faults; 16 candidates takes 30--55\% for 1--2.5\%. The
  sampler is a partial Fisher--Yates over the candidate array and needs no extra memory.
\end{itemize}

\section{Noise floor}
\label{sec:noise}
The linear model was refit from three extra datasets built with different sampling seeds and replayed on
the same test and held-out cells. \IfFileExists{tables/noise.tex}{The faults/Clock geo-mean of the global
K+K+ and F+K models for each data seed is
\begin{center}\small\begin{tabular}{lll}
\toprule
feature group & faults/Clock per data seed & range \\
\midrule
F+K & 0.9452 / 0.9433 / 0.9444 / 0.9459 & 0.0026 \\
K+K+ & 0.8942 / 0.8942 / 0.8945 / 0.8948 & 0.0005 \\
\bottomrule
\end{tabular}
\end{center}
Data sampling moves the result by at most 0.0005 (kernel signals) and 0.0026 (with oracle signals), under 0.3\%, so the
differences between models reported above (1--14\%) are not artefacts of which training rows were drawn. It does not measure
the other sources of variance (the choice of streams, which the bootstrap intervals cover).}{(Result not yet generated.)}

\section{Limits and what is next}
\begin{itemize}
\item The real Redis/SQLite traces and the Twitter replays are small or lack read/write marks; the Twitter streams are
  hash-sampled to 1\,024 keys, which preserves each key's request sequence but not the original working-set size.
\item Neural models were not scored on 1M+ reference streams, so the patbench and appbench groups compare trees only.
\item The kernel runs of patbench and appbench use 50\,000 accesses, not the full 1M, for time.
\item Selection cost is the kernel's selection-ticks counter; a full-eviction-path timer and a no-op policy arm to
  calibrate it are not done.
\item Next: use real kernel runs to decide whether the aging-bit CNN is worth a kernel port (it passes the 2\%
  bar on kv at high cost), and add a heap/allocator-churn program.
\end{itemize}

\section*{Reproducing}
\small
\begin{verbatim}
python3 tools/ml2/zoo_trees.py fit | select | test      # tree + neural zoo (zoo_nn.py fit first)
python3 tools/ml2/zoo_report.py                          # zoo vs linear vs baselines
python3 tools/ml2/exp_transfer.py run [--set real|app]   # transfer streams
python3 tools/ml2/kernel_eval.py --only lzw,kv-T --fracs 0.1  [--models lzw@32]
python3 tools/ml2/kernel_vs_sim.py --only lzw,kv-T,pat,app --tag _w15
python3 tools/ml2/report_w15.py                          # tables and figures of this report
\end{verbatim}
\end{document}
